% WebDAQ manual.
%
% Build with `make` in this directory. The same material is published as a
% website from doc/ (Sphinx); this is the version to print, to hand to a new
% shift crew, or to read on a machine with no network.
\documentclass[11pt,a4paper]{report}

\usepackage[T1]{fontenc}
\usepackage[utf8]{inputenc}
\usepackage[british]{babel}
\usepackage{lmodern}
\usepackage{microtype}
\usepackage[margin=2.6cm]{geometry}
\usepackage{graphicx}
\usepackage{booktabs}
\usepackage{longtable}
\usepackage{array}
\usepackage{enumitem}
\usepackage{fancyvrb}
\usepackage{xcolor}
\usepackage{titlesec}
\usepackage[hidelinks]{hyperref}

\hypersetup{
  pdftitle={WebDAQ — data acquisition for the LUNA experiment},
  pdfauthor={A. Compagnucci, R. Gesuè, J. Skowroński},
  pdfsubject={Operator and technical manual, WebDAQ 5.0},
}

% Monospace for everything the operator types or reads on a screen.
\newcommand{\cmd}[1]{\texttt{#1}}
\newcommand{\ui}[1]{\textbf{#1}}          % a control on the screen
\newcommand{\file}[1]{\texttt{#1}}

\DefineVerbatimEnvironment{shell}{Verbatim}{%
  frame=leftline, framesep=6pt, rulecolor=\color{gray!60},
  fontsize=\small, xleftmargin=4pt}

\newcolumntype{L}[1]{>{\raggedright\arraybackslash}p{#1}}

\titleformat{\chapter}[display]
  {\normalfont\Large\bfseries}{\chaptertitlename\ \thechapter}{8pt}{\LARGE}
\setlist{itemsep=2pt, parsep=0pt, topsep=4pt}

\begin{document}

\begin{titlepage}
\centering
\vspace*{3cm}
{\LARGE\bfseries WebDAQ\par}
\vspace{6mm}
{\Large Data acquisition for the LUNA experiment\par}
\vspace{10mm}
{\large Operator and technical manual\par}
\vspace{2mm}
{\large Version 5.0\par}
\vfill
\includegraphics[width=0.92\textwidth]{../doc/imgs/dashboard.png}
\vfill
{\normalsize A.~Compagnucci, R.~Gesuè, J.~Skowroński\par}
\vspace{2mm}
{\small INFN Laboratori Nazionali del Gran Sasso and Università di Padova\par}
\vspace{6mm}
{\small \today\par}
\end{titlepage}

\tableofcontents

% ============================================================================
\chapter{What WebDAQ is}

WebDAQ reads the CAEN digitizers of the LUNA setup, writes the data, fills the
spectra while they fill, records what the beam was doing, and says something when
a board stops. It is used from a browser, by the person on shift, and it is meant
to be learnt in an evening.

\section{The data path}

\begin{center}
\begin{tabular}{@{}l@{\hspace{6pt}}c@{\hspace{6pt}}l@{}}
digitizers & $\rightarrow$ & CaenDAQ, inside the server process \\
           &               & $\rightarrow$ \file{run\_<N>\_0000.caendat} on disk \\
           &               & $\rightarrow$ spectra, rates, regions of interest \\
           &               & $\rightarrow$ Graphite, for the slow-control archive \\
\end{tabular}
\end{center}

Acquisition runs \emph{inside} the DAQ server. The boards are opened, read,
decoded and written by CaenDAQ, a C\texttt{++} library bound into Python, so
pressing \ui{Start} is the whole operation: there is no container to launch, no
state machine to drive and no second process to keep alive. The spectra you watch
are the histograms CaenDAQ is already accumulating while it decodes; the server
takes a snapshot of them when the page asks.

Versions before~4.0 ran acquisition as XDAQ inside a Docker container, with a
separate spy server on port~6060. Any procedure that mentions a container,
port~6060 or a \file{topology.xml} predates~4.0 and does not apply.

\section{The three pieces}

\begin{itemize}
  \item \textbf{The DAQ server} --- Flask and waitress on port~5001. It owns the
        boards, the run, the files and the configuration.
  \item \textbf{The web interface} --- Next.js on port~3000. It holds no state of
        its own worth protecting; everything it shows comes from the server.
  \item \textbf{RUReader} --- the offline converter from \file{.caendat} to ROOT,
        run on request from the Logbook.
\end{itemize}

\section{The working directory is the experiment}

Everything the server reads and writes is relative to the directory it was started
in: \file{conf/} (board configurations, the histogram dashboard, credentials),
\file{calib/}, \file{data/} and the \file{app.db} database. Two working
directories are two independent experiments that happen to share one
installation.

This is worth knowing when something looks \emph{missing}. A dashboard with no
histograms, or an empty board list, usually means the server was started somewhere
other than where the campaign was configured. That is a different setup, not a
broken one, and the \ui{DAQ Server} panel names the directory it is running in.

% ============================================================================
\chapter{Installing and starting}

\section{What the machine needs}

Linux for a DAQ machine, macOS for development. For real hardware, the CAEN
libraries (CAENVMElib, CAENComm, CAENDigitizer); without them WebDAQ still runs
in test mode with simulated boards. Everything else --- Python, ROOT, the
acquisition module --- lives in a conda environment called \cmd{luna}, created by
the installer.

\section{Installing}

\begin{shell}
git clone https://github.com/skowrons94/WebDAQ.git
cd WebDAQ
./install.sh
\end{shell}

The installer is idempotent: every step checks whether the work is already done,
so running it again after an update is safe. It installs the build tools, checks
out the CaenDAQ and RUReader submodules, installs Miniforge if conda is missing,
creates the \cmd{luna} environment, builds RUReader and the \cmd{caendaq} module,
writes \file{frontend/.env}, builds the interface, and adds a \cmd{LunaDAQ}
launcher to \file{\textasciitilde/.bashrc}.

\paragraph{The API address is compiled in.} \file{frontend/.env} holds
\cmd{NEXT\_PUBLIC\_API\_URL}, and Next.js bakes it into the bundle at build time.
If the server is to be reached from other machines, set it to an address those
machines can resolve --- not \cmd{127.0.0.1} --- and rebuild with
\cmd{npm run build} in \file{frontend/}. Forgetting the rebuild is the most common
deployment mistake: the interface loads, then fails every request.

\section{Starting it}

\begin{shell}
LunaDAQ            # start the web interface (freeing its port first)
LunaDAQ status     # what is holding each port
LunaDAQ stop       # stop both parts, politely then firmly
LunaDAQ restart    # stop, then start the interface
LunaDAQ backend    # run the DAQ server in this terminal, for debugging
\end{shell}

Open \url{http://localhost:3000} and log in. The \ui{DAQ Server} badge at the top
right starts the server itself: choose the working directory for the campaign and
press \ui{Start}. The badge is also on the login page, so the server can be
started before anyone logs in.

\cmd{LunaDAQ stop} only stops processes that look like ours. An unrelated program
on port~3000 is reported and left alone, which is why it is a better answer to a
busy port than killing whatever holds it.

\section{Running it in tmux}

\label{sec:tmux}
The web interface is a long-lived process: started from a plain SSH session it
dies when the connection drops. Start it inside \textbf{tmux}, a session that
keeps running on the machine whether or not anyone is attached to it.

\begin{shell}
tmux new -s daq        # create a session called "daq" and attach to it
LunaDAQ                # start the interface inside it
\end{shell}

Then detach and leave it running: press \textbf{Ctrl-b}, release both keys, and
press \textbf{d}. You are back in your own shell and the session is still
running.

\begin{center}
\begin{tabular}{@{}L{0.46\textwidth}L{0.46\textwidth}@{}}
\toprule
\textbf{What you want} & \textbf{How} \\
\midrule
List the sessions on the machine & \cmd{tmux ls} \\
Attach again & \cmd{tmux attach -t daq} \\
Detach, leaving it running & \textbf{Ctrl-b} then \textbf{d} \\
Take over a session someone left attached & \cmd{tmux attach -d -t daq} \\
Scroll back through the output & \textbf{Ctrl-b} then \textbf{[}, arrows or PageUp; \textbf{q} to leave \\
A second window in the session & \textbf{Ctrl-b} then \textbf{c}; switch with \textbf{Ctrl-b} then \textbf{0}, \textbf{1}, \dots \\
Close a window & type \cmd{exit} in it \\
Stop everything & attach, \cmd{LunaDAQ stop}, then \cmd{exit} \\
\bottomrule
\end{tabular}
\end{center}

Every tmux command is the prefix \textbf{Ctrl-b} followed by one key; that is the
only thing worth memorising. A detached session survives a dropped SSH
connection, but not a reboot of the DAQ machine --- after a reboot, start a new
session.

\section{Without hardware}

\begin{shell}
TEST_FLAG=True LunaDAQ backend
\end{shell}

or tick \ui{Test Mode} in the \ui{DAQ Server} panel before starting. Test mode
fabricates CAEN-style data for simulated boards and simulates the picoammeter, so
the whole chain works on a laptop. Set the variable or leave it unset;
\textbf{never set it to \cmd{False}} --- most of the code reads any value at all
as ``test mode''.

% ============================================================================
\chapter{Taking a run}

\section{Before the beam}

\begin{enumerate}
  \item \ui{Settings $\rightarrow$ Boards}: every board \emph{Connected}. Use
        \ui{Scan for boards} rather than typing settings from memory; the scan
        fills the add form and leaves the ID and the firmware choice to you.
  \item \ui{DAQ $\rightarrow$ Configuration}: thresholds and shaping for the
        measurement. For several boards, check \ui{Synchronization}
        (chapter~\ref{ch:boards}).
  \item \ui{Settings $\rightarrow$ Current Module}: where the beam current comes
        from --- a TetrAMM, an RBD~9103, or a value already monitored in Graphite.
  \item \ui{Stats}: the accelerator values this campaign should record. Their
        name and unit become the column headings of every run's
        \file{stats.csv}.
  \item \ui{Settings $\rightarrow$ Notifications}: at least one destination and
        the alerts worth having (chapter~\ref{ch:watching}).
\end{enumerate}

\section{Starting}

On \ui{Dashboard $\rightarrow$ Overview}, in \ui{Run Setup \& Control}, fill in
\ui{Run Name}, \ui{Run Type}, \ui{TV (kV)} and \ui{PV (kV)} --- all four are
pre-filled from the last run --- and press \ui{Start Run}.

What happens, in order: the run number and the saving options are stored; the
beam-current recording and the statistics are pointed at this run; the boards are
handed to CaenDAQ, configured and started; the run state is set, which starts the
charge integration and the board monitor; and the run's metadata row is written
with the hardware and software provenance read back from the open boards.

Two refusals are worth knowing in advance. A run will not start without a run
number, and a run directory that already holds data is never silently reused ---
you are asked first.

A run also starts when something \emph{next to} it fails. If the current module
or the statistics cannot be started, the run starts anyway and the message names
what is missing: \emph{``Run 1276 is acquiring, but beam current is not being
recorded''}. A run without its current log is worth more than no run.

\section{While it runs}

\begin{center}
\begin{tabular}{@{}L{0.30\textwidth}L{0.62\textwidth}@{}}
\toprule
\textbf{Where} & \textbf{What it answers} \\
\midrule
Overview, board cards & Is every board still \emph{Running}? \\
\ui{Board Health} & What has each board actually read, written and lost? \\
\ui{Data Rates} & Counts per second, per channel, for one board. \\
\ui{Histograms} & Are the spectra filling where they should? \\
\ui{Beam Control} & Current on target, charge this run, lifetime charge. \\
Header & Run state, elapsed time, and the alert bell. \\
\bottomrule
\end{tabular}
\end{center}

Changing settings mid-run is covered in section~\ref{sec:duringrun}.

\section{Stopping}

\ui{Stop Run} stops the monitor, the spy snapshots and the acquisition, hands the
boards back, writes the end time, advances the run number (when saving) and sets
the run state to stopped --- which stores the accumulated charge, closes
\file{current.txt} and writes \file{roi.json}. Auto-managed Grafana rules are
silenced again.

\section{Afterwards}

Open the run in the \ui{Logbook}: notes in Markdown, the beam current it was taken
with, the monitored metrics beside it, the register dump of every board, and
conversion to ROOT. \ui{Logbook $\rightarrow$ ELOG $\rightarrow$ New entry
$\rightarrow$ Write up a run} drafts a logbook entry from the run's own record.

% ============================================================================
\chapter{The screens}

The navigation is the same everywhere: \ui{Dashboard}, \ui{Logbook}, \ui{Stats},
\ui{DAQ}, \ui{Tuner}, \ui{Troubleshoot}, \ui{Alerts}, \ui{Settings}. The header
carries the run state with an elapsed clock, the \ui{DAQ Server} badge, the alert
bell, and a \ui{Search} palette that jumps to any page or tab.

Logging in is per browser and the session does not expire on its own. If it is
cleared, the notice says plainly that the run is unaffected: the DAQ server keeps
taking data with nobody logged in.

\section{Dashboard}

Six tabs: \ui{Overview}, \ui{Board Health}, \ui{Data Rates}, \ui{Histograms},
\ui{Waveforms}, \ui{PSD}. The last three can be switched off in
\ui{Settings $\rightarrow$ Appearance}.

\subsection{Overview}

\ui{System status} is the band of cards at the top; drag the handle beneath it to
show one row more or less.

\begin{center}
\begin{tabular}{@{}L{0.24\textwidth}L{0.68\textwidth}@{}}
\toprule
\textbf{Card} & \textbf{What it tells you} \\
\midrule
\ui{Run Status} & Run number, running or stopped, how long ago it started, the
file write rate, and a \ui{Save} switch. While stopped, the run number is a
button that opens \emph{Next run}. \\
\ui{DAQ Status} & Data saving, waveforms, current, stats --- each with a tick or a
cross, so you can see what this run will actually record. \\
\ui{Picoammeter} & For a real device only: connected or not, whether it is
sampling, its address and range. A monitored Graphite value has no device card. \\
\ui{Board \emph{n}} & \emph{Ready}, \emph{Running}, \emph{Board Failure} or
\emph{Disconnected}, with a waveform switch that is locked during a run. \\
\ui{ROIs}, \ui{Metrics} & Live region-of-interest integrals, and the Graphite
values selected on the Stats page. \\
\bottomrule
\end{tabular}
\end{center}

\ui{Acquisition setup} shows what this run will do --- saving, the file size
limit, the auto-restart delay --- and \ui{Adjust} opens the dialog that sets them.
\ui{Beam Control} plots the current with the charge beside it and can be moved to
one of four places on the page.

\subsection{Board Health}

One card per board with the readout counters CaenDAQ keeps for the run: data
blocks read, events decoded, bytes read and bytes sent to file, each with its
rate, and then the three numbers that matter when something is wrong.

\begin{center}
\begin{tabular}{@{}L{0.22\textwidth}L{0.70\textwidth}@{}}
\toprule
\textbf{Counter} & \textbf{Read it as} \\
\midrule
\ui{FAIL blocks} & The board could not sustain the readout for that block --- a
full buffer or a lost link. Data from that point may be incomplete. \\
\ui{Dropped} & Blocks the write queue refused because it was full. \textbf{This is
lost data} and should always be zero. \\
\ui{Read errors} & CAEN read errors on the link. \\
\bottomrule
\end{tabular}
\end{center}

The badge is the quick answer: \emph{Reading}, \emph{No data for n~s}, \emph{FAIL
flag}, \emph{No answer} (the last request failed, so the numbers are the last
known ones) or \emph{Not reported} (this CaenDAQ build does not publish that
counter). Every counter belongs to the run: they start at zero and are gone when
it ends, which is why the tab says so between runs rather than showing zeros.

\subsection{Data Rates}

A tile per board, four tiles for the selected board --- events, pile-ups, lost
events, saturations --- and the write rate. Pile-up, lost and saturation turn
amber when they are not zero.

The \ui{Plot} control sets the sampling cadence, and it is not a display setting:
one CaenDAQ tick samples the counters, differences them and publishes them, so the
same number sets the refresh rate, the averaging window and the resolution of what
reaches Graphite. \ui{Rate shown} switches between the four rates without losing
the trace, because all four are recorded.

\subsection{Histograms, Waveforms and PSD}

The histogram dashboard belongs to the experiment, not to your browser: the
layout, the regions of interest, the zoom and the display settings live on the DAQ
server in \file{conf/histograms.json}, so every screen in the control room shows
the same thing and a restart changes nothing. Chapter~\ref{ch:spectra} covers it.

\section{Logbook}

Two logbooks: the run metadata WebDAQ records itself, and the collaboration's
ELOG. The run table is searchable and filterable, the flag and the notes are
editable in place, and \ui{Export CSV} writes out exactly the rows you filtered
to. A run number opens the run's own page with six tabs: \ui{Overview},
\ui{Notes}, \ui{Beam Current}, \ui{Stats}, \ui{Boards} and \ui{Convert}
(chapter~\ref{ch:rundata}).

\section{Stats}

The Graphite values this experiment watches, each with its latest reading, a
sparkline and a timestamp. The light at the top says whether Graphite is
answering at all, so a dead server never looks like a metric with no data.
\ui{Add metric} browses or searches the metric tree. The \textbf{name and unit}
you give a metric become the column heading in every run's \file{stats.csv}.

The \ui{Graphite server} card holds the host, the port of the render API, and the
\textbf{metric prefix} for this experiment. Give each campaign its own subtree ---
\cmd{ancillary.rates.12c12c} --- and no two campaigns share a series.

\section{DAQ}

\ui{Configuration} is the full register set of one board, grouped by function, in
the order the signal is processed; \ui{Advanced} shows every register and
\ui{Binary} shows bit fields. \textbf{Writes on this tab reach the board
immediately and are not blocked during a run} --- it is the same dashboard used to
set a board up, so it is not the place to experiment mid-measurement. Use the
Tuner, which knows what is safe.

\ui{Synchronization} is the daisy chain (chapter~\ref{ch:boards}).
\ui{Calibration} holds two numbers per channel for the energy axis.
\ui{Hardware Info} is what the boards report about themselves --- model, serial
number, firmware, licence, and the acquisition registers as programmed. The CAEN
API answers only while the digitizers are open, so that page is populated during a
run.

\section{Tuner}

The place to change thresholds and shaping with beam on target. It starts
acquisition with \textbf{data saving switched off}, so nothing lands in a run
directory, and a badge reminds you that the Tuner is holding the run.

\label{sec:duringrun}
The \ui{Online} switch decides where a change goes. With it on during a run, only
registers on the firmware's safe list are written to the board; every field is
marked \emph{live} or \emph{next run} \emph{before} you touch it. The save to the
configuration always happens, and a write the board refuses is reported as
\emph{saved, not sent to the board} with the reason --- so the configuration and
the hardware cannot drift apart.

\section{Troubleshoot}

\ui{Recovery} at the top, \ui{History} below. Chapter~\ref{ch:trouble}.

\section{Alerts}

The Grafana alert rules, managed without leaving run control. Marking a rule
\textbf{auto-managed} makes it active while a run is going and silent between
runs, which is the answer to a ``beam current too low'' rule that is right during
a run and pure noise while you are setting up.

\section{Settings}

Six sections: \ui{Boards}, \ui{Current Module}, \ui{Notifications}, \ui{ELOG},
\ui{Grafana} and \ui{Appearance}. Everything except Appearance is stored on the
DAQ server, so it is the same for everyone.

\paragraph{Credentials.} \file{conf/telegram\_settings.json},
\file{conf/zulip\_settings.json}, \file{conf/elog\_settings.json} and
\file{conf/grafana\_settings.json} hold secrets in plain text on the DAQ machine.
The interface masks them once saved and never shows them again, but the files are
readable by anyone with an account on that machine --- so do not copy a working
directory between sites without emptying them.

% ============================================================================
\chapter{Boards and synchronisation}
\label{ch:boards}

\section{Adding a board}

\ui{Settings $\rightarrow$ Boards $\rightarrow$ Scan for boards} probes the links
--- USB, optical (CONET), A4818, and a VME crate through a bridge --- and lists
the digitizers that answer with their model, serial number, channel count and
firmware. \ui{Use these settings} fills the add form rather than adding the board,
because the board ID and the DPP firmware are judgements, not measurements: the
firmware is only \emph{inferred} from the code the board reports.

Adding a board reads its register dump out of the hardware into
\file{conf/<model>\_<id>.json} and creates a unit calibration in \file{calib/}.
That JSON file is the source of truth from then on: the Tuner, the CAEN
dashboard, the waveform switch and CaenDAQ itself all read it, and it is copied
into every run directory.

\section{What a synchronised run is}

Synchronisation is never set by WebDAQ. It comes from the boards' own registers,
and the server reads it back rather than remembering what it was told. A board
whose start mode (Acquisition Control \cmd{0x8100}, bits [1:0]) is anything other
than \emph{software} is \emph{armed} rather than started, and waits.

CaenDAQ then starts the run like this: align the clock chain, arm every
synchronised board, and only then fire one software trigger on the \textbf{master}
--- the synchronised board whose register id is~0. That pulse starts the master
and leaves its TRG-OUT, enters the next board's TRG-IN, and walks the chain, so
every board shares one time origin.

\section{What each position needs}

\begin{center}
\begin{tabular}{@{}L{0.26\textwidth}L{0.66\textwidth}@{}}
\toprule
\textbf{Board} & \textbf{Needs} \\
\midrule
The master & \emph{On first trigger} mode, and its software trigger routed to
TRG-OUT (\cmd{0x8110} bit~31) when any board follows it. It does \textbf{not} need
TRG-IN $\rightarrow$ TRG-OUT: no trigger arrives on its TRG-IN, it makes its
own. \\
A board in the middle & TRG-IN $\rightarrow$ TRG-OUT (\cmd{0x8110} bit~30), so the
pulse walks on, and TRG-OUT carrying the trigger rather than a probe
(\cmd{0x811C} bits [17:16] $=0$). \\
The last board & Nothing on its TRG-OUT. The chain ends there. \\
\bottomrule
\end{tabular}
\end{center}

\ui{DAQ $\rightarrow$ Configuration $\rightarrow$ Synchronization} checks all of
this for you and marks each switch \emph{needed --- on}, \emph{needed --- off} or
\emph{not needed here} for that board's position. Anything that would stop the
chain starting is listed against the board that causes it, in words rather than
register numbers. Getting this wrong is expensive and silent --- the run starts,
one board never sees a trigger, and its data is short --- which is why the check
exists.

The master is identified from the register ids the hardware reports. Those are
readable only while the boards are open, so between runs the panel says the master
is \emph{assumed} from the configuration.

\section{Clock as well as start}

Synchronising the \emph{start} aligns the boards at $t=0$; sharing a \emph{clock}
keeps them aligned. On their own oscillators the boards drift apart over a long
run. Leave the first board on \emph{Internal}, feed its CLK-OUT to the next
board's CLK-IN, and set every downstream board to \emph{External CLK-IN}
(\cmd{0x8100} bit~[6]; desktop and NIM digitizers only --- in a VME crate the
clock comes from the crate). CaenDAQ re-aligns the chain at every run start and
reports a board whose PLL did not stay locked, rather than refusing to run.

% ============================================================================
\chapter{Spectra and regions of interest}
\label{ch:spectra}

The spectra are the ones CaenDAQ accumulates while decoding; the server snapshots
them when a page asks. Nothing is transported twice and there is no monitoring
process to fall behind.

\textbf{The dashboard belongs to the experiment.} Which spectra are shown, their
labels, their regions of interest, their zoom and the display settings live in
\file{conf/histograms.json} on the DAQ server. A second screen in the control room
shows the same dashboard, and a restart changes nothing.

\section{Regions of interest}

A region is a name, a range and a colour, defined per spectrum. Integrals are
shown live on the Overview and on the spectrum itself. At every run stop, every
defined region --- visible or not, enabled or not --- is written into the run
directory as \file{roi.json}, with its bounds and its counts, so a run carries the
regions it was analysed with.

\section{An empty spectrum}

An empty histogram says why in its own title: \emph{No data yet --- start a run},
\emph{This board is not part of the current run}, \emph{No waveform --- switch
waveforms on for this board}, or \emph{No counts yet on channel n}. Read the title
first; then check on \ui{Board Health} whether the board is reading data blocks at
all, and only then the channel enable and the threshold.

% ============================================================================
\chapter{Beam current and charge}

\section{Two charges, and why they differ}

\ui{This run} integrates for exactly the length of the run and is stored in the
run's metadata. \ui{Lifetime total} never stops, so it can serve as a target
history. A figure that does not move between runs is correct behaviour.

The run figure follows the DAQ's own run state rather than the browser: a closed
tab, an auto-restart or a run that saves nothing all leave it correct. A dot next
to \emph{This run} says whether it is integrating.

\section{Three sources}

\begin{center}
\begin{tabular}{@{}L{0.22\textwidth}L{0.70\textwidth}@{}}
\toprule
\textbf{Module} & \textbf{What it is} \\
\midrule
TetrAMM & Four-channel picoammeter over TCP. One channel is the charge channel. \\
RBD 9103 & Single-channel picoammeter over a serial port. \\
Monitored value & No hardware: the beam current is read from a value already
published to Graphite --- the accelerator's own readout --- scaled into
microamps. It plots, logs and integrates exactly like a picoammeter. \\
\bottomrule
\end{tabular}
\end{center}

\section{Freshness}

A picoammeter answers in milliseconds, so a reading older than 30~s means it has
stopped. A monitored value is different: it is published on its own cadence,
Carbon flushes on its own schedule, and the render API will not serve the bucket
it is still filling --- so a value written every 10~s legitimately reads back tens
of seconds old. WebDAQ measures the cadence and the lag it actually sees and
derives the freshness horizon from them, which is why a healthy slow metric is no
longer reported as a dead readout.

Non-finite and negative readings are refused rather than integrated, and the
length of gap that may be integrated across is bounded --- so a Graphite outage
cannot add charge that was never delivered.

% ============================================================================
\chapter{Watching the experiment}
\label{ch:watching}

Four different things watch a run, and they are easy to confuse.

\begin{center}
\begin{tabular}{@{}L{0.24\textwidth}L{0.68\textwidth}@{}}
\toprule
\textbf{Tool} & \textbf{What it does here} \\
\midrule
Graphite & Stores every number anyone measures: beam current, terminal voltage,
board rates. WebDAQ writes to it and reads from it. \\
Grafana & Draws those numbers on the control-room screen and raises its own
alerts. \\
The Stats page & Shows the current value of the metrics \emph{you} selected, and
copies them into each run's \file{stats.csv}. \\
Board Health & The readout counters of each board, live for the run only. \\
\bottomrule
\end{tabular}
\end{center}

\section{Alerts are rules}

An alert is a rule: what to watch, the threshold, how long the condition has to
hold, and which destinations carry it. Rules live in \file{conf/alerts.json} and
are edited in \ui{Settings $\rightarrow$ Notifications}.

\begin{center}
\begin{tabular}{@{}L{0.30\textwidth}L{0.62\textwidth}@{}}
\toprule
\textbf{Alert} & \textbf{Fires when} \\
\midrule
Board failure & A board sets its FAIL flag during a run. \\
Board stopped producing data & A board reads no new data block for $n$ seconds
during a run, whether or not it raised the FAIL flag. \\
Beam current & The current stays below (or above) a threshold for $n$ seconds. \\
Graphite metric & A monitored value stays outside its range. One rule per
metric. \\
\bottomrule
\end{tabular}
\end{center}

Destinations are \textbf{Telegram} (a bot posting to a chat or group) and
\textbf{Zulip} (a bot posting to one stream and topic). Each rule names its own,
so a board failure can wake the whole shift while a current dip reaches one group.

\section{Why the messages are worth reading}

\begin{itemize}
  \item \textbf{One message per episode.} A rule fires when its condition starts
        and stays quiet until it clears, then says it cleared.
  \item \textbf{Unknown is not wrong.} An unreachable Graphite, or a monitor that
        stopped, leaves every rule exactly as it was --- a hole in the monitoring
        never invents an alert and never silently clears a real one. Such a rule
        is marked \emph{nothing to measure}, so it is not mistaken for a happy
        one.
  \item \textbf{A flap is one message.} The same subject is not announced twice
        within a minute.
  \item \textbf{An alert that reached nobody says so.} The bell turns amber and
        the history row is badged \emph{reached nobody}.
\end{itemize}

A rule only works while the server is \emph{watching}; the panel says so and
offers \ui{Start watching} when it is not. A board failure is recorded in the
history \textbf{even when no rule is configured for it}.

\section{Auto-restart}

With auto-restart on, a board failure stops the run, flags it \emph{bad} with a
note, and starts a new one after the delay. The new run continues the failed one:
its own \file{current.txt}, charge and \file{stats.csv}, and the target, voltages
and run type carried over. The intent is that a night shift loses one run rather
than the rest of the night.

\section{What happened while nobody was looking}

Every alert and every recovery is written on the server
(\file{conf/alert\_log.json}, the most recent 500 events, kept across restarts).
The \textbf{bell} carries the unread count; \ui{Troubleshoot} is the full history,
with the time, what it was about, which run it belonged to, and which destinations
received it. This is the page to open after a night shift.

% ============================================================================
\chapter{Run data and conversion}
\label{ch:rundata}

\section{What a run leaves behind}

Everything goes into \file{data/run<N>/}, and only when saving is on.

\begin{center}
\begin{longtable}{@{}L{0.30\textwidth}L{0.62\textwidth}@{}}
\toprule
\textbf{File} & \textbf{What it is} \\
\midrule
\endhead
\file{run\_<N>\_0000.caendat} & The data: one unified stream for every board, with
the same header layout the XDAQ-era files had. Further parts appear when a size
limit is set. \\
\file{<model>\_<id>.json} & The board's complete register dump, copied from
\file{conf/} at run start. This is what makes the run reproducible. \\
\file{current.txt} & The beam current, one row per sample, time relative to the
start of the log. \\
\file{stats.csv} & The monitored accelerator values, one row per sampling tick,
with the aliases and units you chose. \\
\file{stats.csv.part1}, \dots & Earlier rows of the same run, kept aside when the
statistics were restarted mid-run. The run report reads them together with
\file{stats.csv}. \\
\file{metadata.json} & The run's record, rewritten after every change, so a
copied directory is self-describing. \\
\file{roi.json} & Every region of interest with its counts, written at stop. \\
\file{run<N>.root} & The converted file, on request. \\
\bottomrule
\end{longtable}
\end{center}

The \file{.caendat} file has no trailer. That is deliberate: a run killed by a
power cut is still readable up to the point where it stopped.

\section{Converting}

\ui{Logbook $\rightarrow$ the run $\rightarrow$ Convert} runs RUReader over every
\file{.caendat} in the directory, in name order so a run split across files keeps
its time order, and writes \file{run<N>.root} next to them. The options offered
are read from the installed converter rather than assumed, so a build that does
not support a flag does not show it. Conversion runs on the server: you can leave
the page.

\section{Checking an interrupted run}

\begin{shell}
cd server
python scripts/check_run_integrity.py data/run1276
\end{shell}

This walks the file structure without converting anything and reports whether the
header, the board definitions and the aggregate chain are consistent --- the first
thing to do after a crash, a full disk or a power cut.

% ============================================================================
\chapter{When something is wrong}
\label{ch:trouble}

\textbf{Start at the Troubleshoot page.} It lists the state of each part of the
system, offers the action that repairs it, and keeps the history of what has
already happened.

\section{The five recovery actions}

\begin{center}
\begin{tabular}{@{}L{0.30\textwidth}L{0.50\textwidth}L{0.12\textwidth}@{}}
\toprule
\textbf{Action} & \textbf{Use it when} & \textbf{During a run} \\
\midrule
Reopen the boards & A board reads \emph{Disconnected} after a link glitch. & No \\
Reset the acquisition & A run failed or the boards are wedged. & No \\
Reconnect the current monitor & The current reads disconnected or stops updating. & Yes \\
Restart the statistics & \file{stats.csv} stopped being written mid-run. & Yes \\
Re-check Graphite & After a Graphite outage. & Yes \\
\bottomrule
\end{tabular}
\end{center}

Nothing here runs by itself. The two that would cost data during a run are refused
with the reason written under the button. Restarting the statistics never
overwrites measurements: the rows already recorded are moved to
\file{stats.csv.part1} first, and if they cannot be moved the action refuses
rather than truncating the file.

\section{Symptoms}

\begin{description}[style=nextline, leftmargin=0pt]
  \item[A board reads \emph{Disconnected}] Nothing opened it, or something else
        holds it. \emph{Reopen the boards} with the run stopped; check that no
        second backend is running (\cmd{LunaDAQ status}); check the cable and the
        link number.
  \item[A board shows \emph{Board Failure}] It set its FAIL flag: it could not
        sustain the readout, and data from that point may be incomplete. Read the
        three problem counters on \ui{Board Health} --- dropped blocks point at
        the disk or the write rate, read errors at the link. The run should
        usually be restarted.
  \item[A board stops producing data without failing] \ui{Board Health} shows
        \emph{No data for n~s}. The FAIL flag was not raised, so suspect the
        trigger, the cabling into TRG-IN, or a chain that never started.
  \item[A synchronised run never starts] \ui{DAQ $\rightarrow$ Synchronization}
        names the board at fault and what it is missing. See
        chapter~\ref{ch:boards}.
  \item[Every metric reads \emph{N/A}] Graphite is not answering --- the light
        says which case it is. \emph{Re-check Graphite} also clears the circuit
        breaker an outage leaves behind.
  \item[An alert never arrived] Check, in this order: is the panel
        \emph{Watching}? Is the rule badged \emph{reaches nobody} (no destination
        on) or \emph{nothing to measure} (nothing to read)? Is the bell amber? The
        history records a board failure even with no rule configured.
  \item[Something looks ``missing''] Almost always the working directory: the
        server was started somewhere else, so it is showing a different
        configuration, not a broken one.
\end{description}

\section{After a crash or a power cut}

Start the server again --- a \cmd{running: true} left behind is cleared with a
warning, because acquisition cannot survive a restart. Check the last run
directory with \cmd{check\_run\_integrity.py}, convert it (the option that skips a
truncated final aggregate is passed automatically when the installed converter
supports it), and read the Troubleshoot history for what the DAQ saw before it
stopped.

\section{The log}

\file{server.log}, in the working directory the server was started in.
\ui{DAQ Server $\rightarrow$ Show logs} shows the last 200~kB of it in the
browser.

% ============================================================================
\appendix

\chapter{Configuration files}

All paths are relative to the working directory.

\begin{center}
\begin{longtable}{@{}L{0.34\textwidth}L{0.58\textwidth}@{}}
\toprule
\textbf{File} & \textbf{Holds} \\
\midrule
\endhead
\file{conf/settings.json} & Run number, data saving, size limit, the board list,
auto-restart and its delay. Written atomically under a lock; an unreadable one is
kept rather than replaced. \\
\file{conf/<model>\_<id>.json} & A board's full register dump. \\
\file{conf/histograms.json} & The histogram dashboard, its ROIs and display
settings. \\
\file{conf/stats.json} & The Graphite render API, the metric prefix, the sampling
cadence and the selected metrics. \\
\file{conf/current.json} & Which current module is active and its settings, plus
the Carbon ingest address and the lifetime charge. \\
\file{conf/alerts.json} & The alert rules. \\
\file{conf/alert\_log.json} & The alert history, most recent 500 events. \\
\file{conf/telegram\_settings.json} & Bot token and chat id. \\
\file{conf/zulip\_settings.json} & Site, bot account, stream and topic. \\
\file{conf/elog\_settings.json} & Logbook URL, shared account, default
attributes. \\
\file{conf/grafana\_settings.json} & Grafana address and service-account token. \\
\file{calib/<model>\_<id>.cal} & Two calibration numbers per channel. \\
\file{app.db} & Users and run metadata. \\
\file{server.log} & The server's log. \\
\bottomrule
\end{longtable}
\end{center}

\chapter{Environment variables}

\begin{center}
\begin{tabular}{@{}L{0.30\textwidth}L{0.62\textwidth}@{}}
\toprule
\textbf{Variable} & \textbf{Effect} \\
\midrule
\cmd{TEST\_FLAG} & Simulated boards and a simulated picoammeter. Set it or leave
it unset; never set it to \cmd{False}. \\
\cmd{DATABASE\_URL} & Where the database lives. Defaults to \file{app.db} in the
working directory. \\
\cmd{SECRET\_KEY}, \cmd{JWT\_SECRET\_KEY} & Flask and token signing keys. Both
have insecure defaults; set them on a machine that is reachable from outside. \\
\cmd{NEXT\_PUBLIC\_API\_URL} & The API address compiled into the interface. \\
\cmd{WEBDAQ\_SKIP\_DB\_UPGRADE} & Skip the schema check at startup. \\
\cmd{WEBDAQ\_FORCE\_RESTART} & Start even over a backend that is taking data. \\
\cmd{RUREADER\_BIN} & Path to the converter, when it is not on \cmd{PATH}. \\
\bottomrule
\end{tabular}
\end{center}

\chapter{Running the tests}

\begin{shell}
cd server
conda activate luna
./tests/run_tests.sh            # everything
./tests/run_tests.sh alerting   # only modules matching test_*alerting*.py
\end{shell}

The suite runs with simulated boards in a scratch working directory, so it never
touches a real \file{conf/}, \file{data/} or database. Tests that drive the HTTP
surface read the directory they are started from, which is why the script creates
one.

\vfill
\noindent\rule{\textwidth}{0.4pt}

\noindent\small For technical support or feature requests:
Jakub Skowroński (\href{mailto:jakub.skowronski@pd.infn.it}{jakub.skowronski@pd.infn.it}),
Alessandro Compagnucci (\href{mailto:alessandro.compagnucci@gssi.it}{alessandro.compagnucci@gssi.it}),
Riccardo Gesuè (\href{mailto:gesue.riccardo@gssi.it}{gesue.riccardo@gssi.it}).

\end{document}
